\begin{ejer}
Demuestra que $D(f)\subseteq D(g)$ si y solo si $f\in\sqrt{(g)}$.
\end{ejer}
\begin{sol}
Gracias a que $D(f)=V(f)^c$ y al \myref{ejer:2.3}:
\[
D(f)\subseteq D(g)
\quad\Longleftrightarrow\quad
V(g)\subseteq V(f)
\quad\Longleftrightarrow\quad
\sqrt{(f)}\subseteq\sqrt{(g)}. \tag{$*$}
\]
Resta probar que $\sqrt{(f)}\subseteq\sqrt{(g)}$ si y solo si $(f)\subseteq\sqrt{(g)}$. En efecto, la necesidad es inmediata de que $(f)\subseteq\sqrt{(f)}$. Para la suficiencia:
\[
(f)\subseteq\sqrt{(g)}%
\quad%
\underset{%
    \mathclap{\substack{%
    \big\uparrow\\\text{monotonía de }\sqrt{\cdot}%
    }}
}{\implies}%
\quad%
\sqrt{(f)} \quad\subseteq\quad \sqrt{\sqrt{(g)}} 
\quad%
\underset{%
    \mathclap{\substack{%
    \big\uparrow\\\text{idempotencia de }\sqrt{\cdot}%
    }}
}{=}%
\quad%
\sqrt{(g)},
\]
Esta equivalencia junto a $(*)$ garantizan:
\[
D(f)\subseteq D(g)
\quad\Longleftrightarrow\quad
(f)\subseteq\sqrt{(g)}
\quad\Longleftrightarrow\quad
f\in\sqrt{(g)}.
\]
\end{sol}